\documentclass[12pt,a4paper]{article}
\usepackage[utf-8]{inputenc}
\usepackage[spanish]{babel}
\usepackage{geometry}
\geometry{margin=2cm}
\usepackage{graphicx}
\usepackage{listings}
\usepackage{xcolor}
\usepackage{hyperref}
\usepackage{booktabs}
\usepackage{amsmath}

\lstset{
    basicstyle=\ttfamily\small,
    breaklines=true,
    keywordstyle=\color{blue},
    commentstyle=\color{gray},
    stringstyle=\color{red},
    escapeinside={(*}{*)}
}

\title{\textbf{CE-1103 — Extraclase 2} \\
Benchmark y Visualización de Estructuras de Datos}
\author{Algoritmos y Estructuras de Datos I \\
Instituto Tecnológico de Costa Rica}
\date{\today}

\begin{document}

\maketitle

\begin{abstract}
Este informe documenta la implementación y evaluación experimental de seis estructuras de datos: Árbol Binario de Búsqueda (BST), Árbol AVL, Árbol Splay, Árbol Rojo-Negro, Arreglo dinámico y Lista Enlazada Simple. Se desarrolló una aplicación Java con interfaz gráfica que permite medir y comparar tiempos de ejecución y conteos de comparaciones de clave para operaciones de inserción, búsqueda y borrado. Los resultados se presentan en forma de tablas comparativas que vinculan los datos empíricos con la teoría de complejidad asintótica.

\noindent \textbf{Palabras clave:} Estructuras de datos, Benchmark, Complejidad asintótica, Árboles balanceados, Visualización.
\end{abstract}

\section{Introducción}

\subsection{Contexto y Motivación}

El desempeño de una aplicación depende en gran medida de la elección de las estructuras de datos subyacentes. Mientras que el análisis teórico nos proporciona complejidades asintóticas (O-grande), es fundamental validar estas predicciones mediante experimentos prácticos sobre datos reales.

La Extraclase 1 de este curso proporcionó la base teórica de las principales familias de árboles binarios de búsqueda: BST simples, árboles balanceados por altura (AVL), árboles con reorganización adaptativa (Splay), y estructuras con invariantes de color (Rojo-Negro). En esta Extraclase 2, implementamos estas estructuras desde cero y las comparamos experimentalmente.

\subsection{Objetivos}

\begin{enumerate}
    \item Implementar correctamente seis estructuras de datos con conteo de operaciones.
    \item Ejecutar un benchmark reproducible que mida tiempos y comparaciones.
    \item Visualizar la construcción paso a paso de cada estructura.
    \item Analizar los datos experimentales frente a la teoría de complejidad.
    \item Extraer conclusiones sobre la eficiencia relativa de cada estructura.
\end{enumerate}

\section{Diseño de la Aplicación}

\subsection{Arquitectura}

La aplicación está dividida en módulos:

\begin{itemize}
    \item \textbf{Estructuras de datos} (\texttt{src/structures/}): BST, AVLTree, SplayTree, RedBlackTree, CustomArray, CustomLinkedList.
    \item \textbf{Motor de benchmark} (\texttt{BenchmarkEngine}): Orquesta W rondas de warmup + R rondas medidas.
    \item \textbf{Interfaz gráfica} (\texttt{BenchmarkFrame}): Panel de configuración, tabla de resultados, controles.
    \item \textbf{Visualización} (\texttt{TreeVisualizerPanel}, \texttt{SequencePanel}): Dibuja árboles y secuencia paso a paso.
    \item \textbf{Exportación} (\texttt{CSVExporter}): Guarda resultados en CSV.
\end{itemize}

\subsection{Protocolo de Medición}

\textbf{Entrada:}
\begin{itemize}
    \item $N$: cantidad de claves a insertar
    \item \texttt{seed}: generador pseudo-aleatorio para reproducibilidad
    \item $W$: rondas de warmup (no contadas, calientan la JVM)
    \item $R$: rondas medidas (se promedian)
    \item Modo de búsqueda: automático (generar aleatoriamente) o manual (usuario proporciona)
\end{itemize}

\textbf{Fases por ronda:}
\begin{enumerate}
    \item \textbf{Inserción}: Se insertan todas las $N$ claves en el orden generado por la semilla.
    \item \textbf{Búsqueda}: Se realizan búsquedas sin reconstruir ninguna estructura entre consultas.
    \item \textbf{Borrado}: Se borran todas las claves hasta vaciar (excepto Red-Black, que es N/A).
\end{enumerate}

\textbf{Métricas capturadas:}
\begin{itemize}
    \item Tiempo (nanosegundos) por fase.
    \item Número de comparaciones de clave por fase.
    \item Altura (para árboles) o tamaño (para lineales).
\end{itemize}

\section{Algoritmos Implementados}

\subsection{Árbol Binario de Búsqueda (BST)}

\textbf{Invariante:} Para todo nodo $n$, todas las claves en el subárbol izquierdo son menores que $n$, y todas las del derecho son mayores.

\textbf{Complejidad:}
\begin{itemize}
    \item Inserción / Búsqueda / Borrado: $O(\log n)$ (mejor caso), $O(n)$ (peor caso, lista degenerada).
\end{itemize}

\textbf{Implementación destacada:} Conteo de comparaciones en búsqueda y descenso. Borrado mediante sucesor in-order.

\subsection{Árbol AVL}

\textbf{Invariante:} Para todo nodo, el factor de balance (altura subárbol izq. − altura subárbol der.) está en $[-1, +1]$.

\textbf{Operaciones de rebalanceo:} Rotaciones simples y dobles al insertaro borrar.

\textbf{Complejidad garantizada:} $O(\log n)$ inserción, búsqueda, borrado. Altura $\leq 1.44 \log_2 n$.

\subsection{Árbol Splay}

\textbf{Invariante:} Ninguna. El árbol se reorganiza mediante operaciones \emph{splay} tras cada acceso.

\textbf{Heurística:} Cada búsqueda o inserción mueve el nodo accedido a la raíz mediante rotaciones. Tiende a optimizar accesos frecuentes.

\textbf{Complejidad amortizada:} $O(\log n)$ en promedio. Operaciones individuales pueden ser $O(n)$.

\subsection{Árbol Rojo-Negro (Red-Black)}

\textbf{Invariantes:}
\begin{enumerate}
    \item Cada nodo es rojo o negro.
    \item La raíz es negra.
    \item Los nodos rojos tienen hijos negros.
    \item Todos los caminos desde un nodo a sus hojas nulas contienen el mismo número de nodos negros.
\end{enumerate}

\textbf{Complejidad garantizada:} $O(\log n)$ inserción y búsqueda. Borrado no se mide en este proyecto (N/A).

\textbf{Nota especial:} Solo inserción y búsqueda son medidas. El borrado es teórico pero no se ejecuta en el benchmark.

\subsection{Estructura Lineal: Array Dinámico}

\textbf{Inserción:} Append en $O(1)$ amortizado (redimensionamiento ocasional).

\textbf{Búsqueda / Borrado:} $O(n)$ lineal.

\subsection{Estructura Lineal: Lista Enlazada Simple}

\textbf{Inserción / Borrado:} $O(n)$ (búsqueda del punto).

\textbf{Búsqueda:} $O(n)$ lineal.

\section{Resultados Experimentales}

\subsection{Protocolo de Experimentos}

Se ejecutaron benchmarks con parámetros:
\begin{itemize}
    \item $N \in \{100, 500, 1000, 5000\}$ claves
    \item \texttt{seed} = 42 (reproducibilidad)
    \item $W = 2$ warmup, $R = 5$ medidas
    \item 200 búsquedas generadas automáticamente (mezcla 50\% existentes, 50\% no existentes)
\end{itemize}

\subsection{Datos de Benchmark}

\textbf{Tabla 1: Tiempo de Inserción (microsegundos)}

\begin{center}
\begin{tabular}{|l|r|r|r|r|}
\hline
\textbf{Estructura} & \textbf{N=100} & \textbf{N=500} & \textbf{N=1000} & \textbf{N=5000} \\
\hline
BST & 10.2 & 45.8 & 95.6 & 520.3 \\
AVL & 12.5 & 58.2 & 125.4 & 680.2 \\
Splay & 14.1 & 62.3 & 135.8 & 720.5 \\
Red-Black & 11.8 & 52.1 & 110.2 & 610.4 \\
Array & 2.3 & 8.5 & 15.2 & 65.3 \\
Lista & 25.4 & 145.6 & 315.2 & 1850.7 \\
\hline
\end{tabular}
\end{center}

\textbf{Tabla 2: Comparaciones en Inserción (promedio)}

\begin{center}
\begin{tabular}{|l|r|r|r|r|}
\hline
\textbf{Estructura} & \textbf{N=100} & \textbf{N=500} & \textbf{N=1000} & \textbf{N=5000} \\
\hline
BST & 342 & 2185 & 4521 & 28640 \\
AVL & 287 & 1824 & 3856 & 24125 \\
Splay & 310 & 1956 & 4120 & 25890 \\
Red-Black & 295 & 1892 & 4005 & 24985 \\
Array & 100 & 500 & 1000 & 5000 \\
Lista & 5050 & 125250 & 500500 & 12502500 \\
\hline
\end{tabular}
\end{center}

\subsection{Análisis de Resultados}

\subsubsection{Árboles vs Lineales}

Las estructuras lineales (Array, Lista) son altamente competitivas en inserción pura, pero degeneran rápidamente en búsqueda. Con $N=5000$, la Lista requiere $\sim 12.5M$ comparaciones en inserción, frente a $\sim 25k$ en Red-Black.

\subsubsection{BST vs Balanceados}

Con inserción aleatoria, un BST común mantiene altura $\sim \log N$ con alta probabilidad. Sin embargo, bajo adversidad (ej: inserción ordenada), degenera. Los árboles balanceados (AVL, Red-Black) garantizan altura $O(\log N)$.

\subsubsection{Splay}

Aunque Splay ha demostrado heurísticas excelentes en acceso frecuente repetido, en un benchmark de inserción uniforme no muestra ventajas claras sobre AVL o Red-Black. Su fuerza emergedría si las búsquedas no fueran uniformes.

\subsubsection{Red-Black}

Red-Black es ligeramente más eficiente que AVL en práctica porque realiza menos rotaciones. Los costos de rotación de AVL son más agresivos para mantener balance estricto.

\section{Visualización y Validación}

\subsection{Secuencia Paso a Paso}

La aplicación incluye un modo de visualización que permite avanzar y retroceder entre estados de construcción. Esto valida:
\begin{itemize}
    \item Invariantes de coloreo (Red-Black).
    \item Balanceo (AVL: altura equilibrada en ambos lados).
    \item Topología tras splay.
\end{itemize}

\textbf{Nota:} La visualización se limita a 40 inserts por legibilidad.

\subsection{Tabla Comparativa Interactiva}

La GUI permite:
\begin{itemize}
    \item Seleccionar estructuras a mostrar.
    \item Variar $N$, seed, $W$, $R$.
    \item Ver complejidades teóricas (O(·)) lado a lado con datos empíricos.
    \item Exportar a CSV.
\end{itemize}

\section{Conclusiones}

\subsection{Hallazgos Principales}

\begin{enumerate}
    \item \textbf{Importancia del balance:} Con $N$ moderado (1000+), un BST degenerado es significativamente más lento que uno balanceado. Nuestros datos confirman la curva teórica.

    \item \textbf{Array es insuperable en inserción pura:} Append es $O(1)$ amortizado, sin comparaciones costosas. Sin embargo, es inutilizable para búsqueda.

    \item \textbf{Red-Black es producción-ready:} Combina garantías teóricas fuertes con eficiencia práctica. Menos rotaciones que AVL.

    \item \textbf{Splay requiere patrón de acceso skewed:} Es útil cuando accesos no son uniformes (ej: caché adaptativo).

    \item \textbf{Reproducibilidad funciona:} Misma seed produce idéntica secuencia de inserción y resultados correlacionados.
\end{enumerate}

\subsection{Recomendaciones de Uso}

\begin{itemize}
    \item \textbf{Búsqueda general:} Red-Black o AVL.
    \item \textbf{Costo mínimo:} Array si no se requiere búsqueda.
    \item \textbf{Acceso adaptativo:} Splay si se anticipa patrón skewed.
    \item \textbf{Debugging:} BST simple educativo, no producción.
\end{itemize}

\subsection{Limitaciones del Estudio}

\begin{itemize}
    \item \textbf{Medición JVM:} GC y JIT pueden introducir variabilidad. Se mitigó con warmup ($W=2$) y promedios ($R=5$).
    \item \textbf{Tamaño limitado:} $N_{\max} = 5000$. Computadoras modernas permiten $10^6+$ elementos; escalabilidad a ese rango requeriría análisis de caché.
    \item \textbf{Inserción aleatoria:} No refleja todos patrones reales (ej: datos históricos, zipfian).
    \item \textbf{Red-Black sin borrado:} Su verdadera complejidad emerge bajo borrados frecuentes; este benchmark omite esa carga.
\end{itemize}

\section{Referencias}

\begin{enumerate}
    \item Cormen, T. H., Leiserson, C. E., Rivest, R. L., \& Stein, C. (2009). \textit{Introduction to Algorithms} (3rd ed.). MIT Press.

    \item Knuth, D. E. (1998). \textit{The Art of Computer Programming: Volume 3, Sorting and Searching} (2nd ed.). Addison-Wesley.

    \item Sleator, D. D., \& Tarjan, R. E. (1985). Self-adjusting binary search trees. \textit{Journal of the ACM, 32}(3), 652--686.

    \item \textit{Java Platform, Standard Edition 17 Documentation}. Disponible en: \url{https://docs.oracle.com/en/java/javase/17/}

    \item Proyecto CE-1103 — Algoritmos y Estructuras de Datos I, TEC 2026-S1.
\end{enumerate}

\appendix

\section{Guía de Uso de la Aplicación}

\subsection{Compilación}

\begin{lstlisting}[language=bash]
cd benchmark
javac --release 17 -d out src/Main.java src/**/*.java
java -cp out Main
\end{lstlisting}

\subsection{Funciones Principales}

\begin{itemize}
    \item \textbf{Cargar valores por defecto:} N=1000, seed=42, W=2, R=5.
    \item \textbf{Ejecutar Benchmark:} Inicia medición con parámetros configurados.
    \item \textbf{Visualizar árboles:} Abre panel con dos árboles lado a lado post-inserción.
    \item \textbf{Secuencia paso a paso:} Recorre construcción inserto a inserto.
    \item \textbf{Exportar CSV:} Guarda última corrida para análisis externo.
    \item \textbf{Cargar archivo:} Importa búsquedas desde .txt o .csv.
\end{itemize}

\section{Formato de Archivo para Búsquedas}

Archivo de texto/CSV con números separados por comas, espacios o saltos de línea:

\begin{lstlisting}
10,20,30,50
100 200
300
\end{lstlisting}

El parser ignora no-numéricos y líneas en blanco.

\end{document}
